\documentclass{Group_Templ}

\usepackage{amsthm}
\usepackage{amssymb}

\begin{document}

\section*{Task 2.4}


\begin{proof}[Augmenting and Directed Paths]
	We want to prove the following statement:
	
	Let $G = (V, E)$ be a bipartite graph with
	$V = A \dot\cup B$, let $D$ be the directed graph
	constructed from $G$, and let $M$ be a matching in $G$.
	
	An augmenting path exists in $G$ with respect to $M$
	if and only if there is a directed path in $D$
	between an exposed vertex in $A$ and an exposed
	vertex in $B$.
	
	\medskip
	\noindent\textbf{Construction of the Directed Graph $D$}
	
	The edges of $G$ are directed according to their
	membership in $M$:
	
	\begin{itemize}
		\item Every edge $e \notin M$ is directed
		from $A$ to $B$.
		\item Every edge $e \in M$ is directed
		from $B$ to $A$.
	\end{itemize}
	
	This orientation ensures that a directed path from
	an exposed vertex in $A$ to an exposed vertex in $B$
	alternates between matching and non-matching edges.
	
	\medskip
	\noindent\textbf{Direction 1:
		Augmenting Path $\Rightarrow$ Directed Path}
	
	Assume that an augmenting path $P$ exists in $G$
	with respect to $M$.
	
	By definition, an augmenting path starts and ends
	at exposed vertices. Furthermore, its edges alternate
	between edges belonging to $M$ and edges not
	belonging to $M$.
	
	Since $G$ is bipartite and an augmenting path starts
	and ends with an edge not belonging to $M$, one
	endpoint belongs to $A$ and the other belongs to $B$.
	
	We traverse the path from $A$ to $B$. According to
	the construction of $D$, all non-matching edges are
	directed from $A$ to $B$, while all matching edges
	are directed from $B$ to $A$.
	
	Therefore, every edge of $P$ follows the
	corresponding direction in $D$.
	
	Thus, $P$ is also a directed path in $D$ between
	an exposed vertex in $A$ and an exposed vertex
	in $B$.
	
	\medskip
	\noindent\textbf{Direction 2:
		Directed Path $\Rightarrow$ Augmenting Path}
	
	Assume that there exists a directed path $P$ in $D$
	from an exposed vertex in $A$ to an exposed
	vertex in $B$.
	
	According to the construction of $D$:
	
	\begin{itemize}
		\item Edges directed from $A$ to $B$ do not
		belong to $M$.
		\item Edges directed from $B$ to $A$ belong
		to $M$.
	\end{itemize}
	
	Since $G$ is bipartite, the path alternates between
	vertices in $A$ and $B$. Therefore, its edges
	alternate between matching and non-matching edges.
	
	Furthermore, both endpoints of $P$ are
	exposed vertices.
	
	Consequently, $P$ satisfies all conditions of
	an augmenting path in $G$.
	
	\medskip
	\noindent\textbf{Conclusion}
	
	Thus, both directions have been proven:
	
	\[
	\boxed{
		\begin{gathered}
			\text{Augmenting path in } G\\
			\iff\\
			\text{Directed path in } D\\
			\text{between exposed vertices in }
			A \text{ and } B
		\end{gathered}
	}
	\]
\end{proof}
\newpage


\section*{Task 2.5}

\noindent\textbf{Initial Matching}

The vertex partitions are:
\[
\begin{aligned}
	A &= \{1,2,3,4,5,6,7\},\\
	B &= \{8,9,10,11,12,13,14\}.
\end{aligned}
\]

The initial matching consists of the red edges:
\[
M_0 = \{(2,8),(5,10),(7,13)\}.
\]

Therefore:
\[
|M_0| = 3.
\]

The exposed vertices are:
\[
\begin{aligned}
	A_{\mathrm{free}} &= \{1,3,4,6\},\\
	B_{\mathrm{free}} &= \{9,11,12,14\}.
\end{aligned}
\]

We direct non-matching edges from $A$ to $B$
and matching edges from $B$ to $A$.

\medskip
\noindent\textbf{Step 0 -- Directed Graph $D_0$}

The directed edges are:
\[
\begin{aligned}
	D_0 = \{&
	1\rightarrow 8,\ 1\rightarrow 10,\
	8\rightarrow 2,\ 2\rightarrow 9,\\
	&3\rightarrow 9,\ 3\rightarrow 12,\
	4\rightarrow 9,\ 4\rightarrow 11,\\
	&10\rightarrow 5,\ 5\rightarrow 14,\
	6\rightarrow 11,\\
	&7\rightarrow 12,\ 13\rightarrow 7
	\}.
\end{aligned}
\]

We identify the following augmenting path:
\[
\boxed{
	P_1: 1\rightarrow 10\rightarrow 5\rightarrow 14
}
\]

We remove the matching edge $(5,10)$ and add
the edges $(1,10)$ and $(5,14)$.

The updated matching is:
\[
\begin{aligned}
	M_1 &= \{(2,8),(1,10),(5,14),(7,13)\},\\
	|M_1| &= 4.
\end{aligned}
\]

\medskip
\noindent\textbf{Step 1 -- Directed Graph $D_1$}

After updating the matching, the directions
of the affected edges change:
\[
\begin{aligned}
	D_1 = \{&
	1\rightarrow 8,\ 10\rightarrow 1,\
	8\rightarrow 2,\ 2\rightarrow 9,\\
	&3\rightarrow 9,\ 3\rightarrow 12,\
	4\rightarrow 9,\ 4\rightarrow 11,\\
	&5\rightarrow 10,\ 14\rightarrow 5,\
	6\rightarrow 11,\\
	&7\rightarrow 12,\ 13\rightarrow 7
	\}.
\end{aligned}
\]

The exposed vertices are now:
\[
\begin{aligned}
	A_{\mathrm{free}} &= \{3,4,6\},\\
	B_{\mathrm{free}} &= \{9,11,12\}.
\end{aligned}
\]

We choose the augmenting path:
\[
\boxed{
	P_2: 3\rightarrow 12
}
\]

Since both endpoints are exposed, this edge
can be added directly to the matching:
\[
\begin{aligned}
	M_2 &= M_1\cup\{(3,12)\},\\
	|M_2| &= 5.
\end{aligned}
\]

\medskip
\noindent\textbf{Step 2 -- Directed Graph $D_2$}

Compared to $D_1$, only the direction of
the edge $(3,12)$ changes:
\[
D_2 =
(D_1\setminus\{3\rightarrow 12\})
\cup\{12\rightarrow 3\}.
\]

The exposed vertices are:
\[
\begin{aligned}
	A_{\mathrm{free}} &= \{4,6\},\\
	B_{\mathrm{free}} &= \{9,11\}.
\end{aligned}
\]

We select another augmenting path:
\[
\boxed{
	P_3: 4\rightarrow 9
}
\]

The updated matching becomes:
\[
\begin{aligned}
	M_3 &= M_2\cup\{(4,9)\},\\
	|M_3| &= 6.
\end{aligned}
\]

\medskip
\noindent\textbf{Step 3 -- Directed Graph $D_3$}

The direction of the edge $(4,9)$ changes:
\[
D_3 =
(D_2\setminus\{4\rightarrow 9\})
\cup\{9\rightarrow 4\}.
\]

The remaining exposed vertices are:
\[
\begin{aligned}
	A_{\mathrm{free}} &= \{6\},\\
	B_{\mathrm{free}} &= \{11\}.
\end{aligned}
\]

We choose the final augmenting path:
\[
\boxed{
	P_4: 6\rightarrow 11
}
\]

We add this edge to the matching:
\[
\begin{aligned}
	M_4 &= M_3\cup\{(6,11)\},\\
	|M_4| &= 7.
\end{aligned}
\]

\medskip
\noindent\textbf{Step 4 -- Final Directed Graph $D_4$}

The last matching edge is reversed:
\[
D_4 =
(D_3\setminus\{6\rightarrow 11\})
\cup\{11\rightarrow 6\}.
\]

There are no exposed vertices remaining.

The final matching is:
\[
\boxed{
	\begin{aligned}
		M_4 = \{&
		(1,10),(2,8),(3,12),\\
		&(4,9),(5,14),(6,11),(7,13)
		\}.
	\end{aligned}
}
\]

\medskip
\noindent\textbf{Conclusion}

Since the graph contains 14 vertices, any matching
can contain at most seven edges.

We have found a matching containing exactly
seven edges:
\[
\boxed{
	|M_4| = 7 = \frac{|V|}{2}
}
\]

Therefore, $M_4$ is a perfect matching and
consequently also a maximum matching.

\end{document}
